\chapter{Equilibrium Benchmarks}
\label{ch:equilibria}

This appendix derives the rational-bidder equilibrium benchmarks that the results chapters compare against. The derivations are standard and are collected here to keep the main text focused on the experimental findings.

\section{Constant Valuations}
\label{sec:eq_constant}

Experiment~1a uses constant valuations $v_i = 1$ with the discrete bid grid $\mathcal{B} = \{0, 0.1, \ldots, 1.0\}$. In a first-price auction with $n$ bidders, a symmetric profile in which every bidder bids $b$ yields each bidder an expected payoff of $(1 - b) / n$, since ties are broken uniformly at random. A unilateral upward deviation to $b + 0.1$ (provided $b + 0.1 \leq 1$) yields a payoff of $1 - (b + 0.1)$. No upward deviation is profitable when
\[
1 - (b + 0.1) \;\leq\; \frac{1 - b}{n},
\]
which after rearrangement gives
\[
b \;\geq\; \frac{0.9\,n - 1}{n - 1}.
\]
For any $b$ at or above this threshold and bounded by the grid, the symmetric profile is a pure-strategy Nash equilibrium. With $n = 2$ the threshold is $0.8$, and with $n = 4$ the threshold is approximately $0.87$.

In a second-price auction with identical valuations, bidding $v_i = 1$ is weakly dominant, since the winner pays the second-highest bid regardless of her own bid. Any profile in which at least one bidder bids $1$ is a Nash equilibrium, and the expected revenue to the seller is $1$. Under revenue equivalence \citep{Vickrey1961} the two formats deliver the same expected revenue, so the first-price equilibrium threshold derivation above pins down the minimum symmetric bid consistent with non-dominated play.

\section{Affiliated Valuations}
\label{sec:eq_affiliated}

Experiments~1b, 2a, and 2b use the linear-affiliation model of \Cref{sec:env}. Each bidder draws a signal $s_i$ uniformly on $\{0, 0.1, \ldots, 1.0\}$, and her valuation is
\[
v_i \;=\; \alpha \, s_i \;+\; \beta \sum_{j \neq i} s_j,
\qquad
\alpha = 1 - \frac{\eta}{2},
\qquad
\beta = \frac{\eta}{2 (n - 1)},
\]
with $\eta \in \{0, 0.5, 1\}$. Signals are drawn independently across bidders and rounds. The symmetric Bayes-Nash equilibrium in both formats features a linear bid function.

\subsection{Bid Functions}

In the second-price auction with independent private values, each bidder has a dominant strategy to bid her expected valuation conditional on the event that she ties for the highest bid. Under the linear model above, the resulting symmetric bid function is
\[
b^{\mathrm{SPA}}(s) \;=\; \Big(\alpha + \frac{n \beta}{2}\Big) s.
\]
The derivation conditions on the highest rival signal equalling the bidder's own signal. Her own signal contributes $\alpha s$, the tying rival contributes $\beta s$, and each of the remaining $n - 2$ rivals contributes $\beta \, \mathbb{E}[s_j \mid s_j \leq s] = \beta s / 2$. Summing across the $n - 1$ rival contributions gives $\beta (1 + (n-2)/2) s = \beta (n/2) s$, which yields the formula above.

In the first-price auction under independent uniform signals, the standard envelope argument produces
\[
b^{\mathrm{FPA}}(s) \;=\; \frac{n - 1}{n} \Big(\alpha + \frac{n \beta}{2}\Big) s.
\]
The bid function is linear in the signal, with a shading factor $(n - 1)/n$ that captures the first-price markdown.

\subsection{Expected Revenue}

Let $\phi = \alpha + n \beta / 2$. Under independent uniform signals on $[0, 1]$, the highest signal has expected value $n / (n + 1)$ and the second-highest signal has expected value $(n - 1) / (n + 1)$. Expected revenue in each format is the bid slope times the relevant order statistic, which gives
\[
R^{\mathrm{SPA}} \;=\; \phi \cdot \frac{n - 1}{n + 1},
\qquad
R^{\mathrm{FPA}} \;=\; \frac{n - 1}{n} \cdot \phi \cdot \frac{n}{n + 1} \;=\; \phi \cdot \frac{n - 1}{n + 1}.
\]
The two expressions coincide, so revenue equivalence holds for every affiliation level $\eta$ in this model. The equivalence holds because signals are drawn independently, which is why the classical \citet{MilgromWeber1982} ranking in favour of second-price formats does not apply here; their ranking requires positively affiliated signals in the strict technical sense. The linear-affiliation function I use creates interdependent valuations but preserves signal independence.

\subsection{Efficient Benchmark}

The expected value of the highest valuation is
\[
\mathbb{E}[v_{(1)}] \;=\; (\alpha - \beta) \, \mathbb{E}[s_{(1:n)}] \;+\; \beta \, \mathbb{E}[S],
\]
where $s_{(1:n)}$ is the maximum of $n$ independent uniform signals and $S = \sum_i s_i$. Substituting $\mathbb{E}[s_{(1:n)}] = n / (n + 1)$ and $\mathbb{E}[S] = n / 2$, the expression simplifies to
\[
\mathbb{E}[v_{(1)}] \;=\; (\alpha - \beta) \cdot \frac{n}{n + 1} \;+\; \beta \cdot \frac{n}{2}.
\]
This is the efficient first-best benchmark. The factorial analysis in Chapter~\ref{ch:qlearning} reports realised revenues rather than welfare, so the efficient benchmark enters the discussion only for comparison with the Price of Anarchy analysis in Chapter~\ref{ch:pacing}.

\section{Numerical Verification}
\label{sec:eq_numerical}

The analytical equilibria above are verified by numerical simulation. The deviation check fixes a grid of signal types and computes the expected payoff from multiplicative deviations around the predicted equilibrium bid. \Cref{tab:bne_deviation} reports the maximal gain from deviating for each $(\eta, n)$ pair; values at numerical zero indicate that the proposed bid is a best response at the reported precision.

\begin{table}[H]
  \centering
  \small
  \caption{Maximum payoff gain from unilateral deviation. Values near zero confirm BNE optimality (200K MC draws per configuration).}
  \label{tab:bne_deviation}
  \begin{tabular}{ccccc}
    \toprule
    $\eta$ & $N$ & Auction & Bid slope & Max Gain \\
    \midrule
    0.0 & 2 & FIRST & 0.5000 & 0.000000 \\
    0.0 & 2 & SECOND & 1.0000 & 0.000000 \\
    0.0 & 3 & FIRST & 0.6667 & 0.000129 \\
    0.0 & 3 & SECOND & 1.0000 & 0.000000 \\
    0.0 & 6 & FIRST & 0.8333 & 0.000002 \\
    0.0 & 6 & SECOND & 1.0000 & 0.000000 \\
    0.5 & 2 & FIRST & 0.5000 & 0.000034 \\
    0.5 & 2 & SECOND & 1.0000 & 0.000000 \\
    0.5 & 3 & FIRST & 0.6250 & 0.000024 \\
    0.5 & 3 & SECOND & 0.9375 & 0.000001 \\
    0.5 & 6 & FIRST & 0.7500 & 0.000000 \\
    0.5 & 6 & SECOND & 0.9000 & 0.000000 \\
    1.0 & 2 & FIRST & 0.5000 & 0.000084 \\
    1.0 & 2 & SECOND & 1.0000 & 0.000000 \\
    1.0 & 3 & FIRST & 0.5833 & 0.000020 \\
    1.0 & 3 & SECOND & 0.8750 & 0.000002 \\
    1.0 & 6 & FIRST & 0.6667 & 0.000003 \\
    1.0 & 6 & SECOND & 0.8000 & 0.000000 \\
    \bottomrule
  \end{tabular}
\end{table}

The revenue check runs Monte Carlo simulations at the equilibrium bids and compares realised mean revenue to the closed-form prediction $R^{\mathrm{BNE}} = \phi \cdot (n - 1)/(n + 1)$ for both formats. \Cref{tab:bne_revenue_match} reports the results.

\begin{table}[H]
  \centering
  \small
  \caption[Revenue formula validation: Analytical vs.\ Monte Carlo (500K draws)]{\textbf{Revenue formula validation: Analytical vs.\ Monte Carlo (500K draws).} Revenue equivalence holds under iid signals.}
  \label{tab:bne_revenue_match}
  \begin{tabular}{cccccc}
    \toprule
    $\eta$ & $N$ & $R^{\text{analytical}}$ & $R^{\text{FPA}}_{\text{MC}}$ & $R^{\text{SPA}}_{\text{MC}}$ & $|R^{\text{FPA}} - R^{\text{SPA}}|$ \\
    \midrule
    0.0 & 2 & 0.3333 & 0.3330 & 0.3337 & 0.0007 \\
    0.0 & 3 & 0.5000 & 0.5000 & 0.4997 & 0.0003 \\
    0.0 & 6 & 0.7143 & 0.7145 & 0.7145 & 0.0000 \\
    0.5 & 2 & 0.3333 & 0.3330 & 0.3337 & 0.0007 \\
    0.5 & 3 & 0.4688 & 0.4688 & 0.4685 & 0.0003 \\
    0.5 & 6 & 0.6429 & 0.6431 & 0.6430 & 0.0000 \\
    1.0 & 2 & 0.3333 & 0.3330 & 0.3337 & 0.0007 \\
    1.0 & 3 & 0.4375 & 0.4375 & 0.4372 & 0.0003 \\
    1.0 & 6 & 0.5714 & 0.5716 & 0.5716 & 0.0000 \\
    \bottomrule
  \end{tabular}
\end{table}

The Monte Carlo means match the closed-form predictions within tight confidence intervals, and the first-price and second-price revenues coincide within sampling error, confirming revenue equivalence.

\section{Pacing Equilibrium}
\label{sec:eq_pacing}

The pacing experiments in Chapter~\ref{ch:pacing} do not admit a closed-form Bayes-Nash benchmark, because budget constraints couple the bidder's choices across rounds. Instead, the relevant equilibrium concept is the first-price pacing equilibrium introduced by \citet{Conitzer2022FPPE}, which is unique under mild regularity conditions and computable via the Eisenberg-Gale convex programme. The welfare benchmark is the liquid welfare optimum defined in \Cref{sec:welfare}, and the efficiency benchmark is the Price of Anarchy bound of $\mathrm{PoA} \leq 2$ derived by \citet{Aggarwal2019} and \citet{Gaitonde2023}. I do not rederive these results here; the reader is referred to the original sources for the full constructions.
